\section{Conclusion}

The Vera C.\ Rubin Observatory will produce science from one of the largest and fastest data
systems ever built for astronomy \cite{rubin-dm,rubin-keynumbers,rubin-alerts}. Operating that
system means being able to observe it---not just the applications, but the storage, the networks,
the facility services, and everything in between.

The platform described here is practical and continues to evolve. It is built from open-source
tools, Kubernetes conventions, and GitOps workflows, and it covers metrics, logs, dashboards, alert
rules, long-term storage, and operational views across a wide range of systems
\cite{lsst-it-k8s-cookbook,prometheus-operator,grafana-loki-alloy}. It is not finished, and it was
never meant to be.

The most important thing is not which tools were chosen. It is that the team built a way to keep
asking better questions. What changed? Where did the latency start? Which system is actually
affected? Who needs to act, and what do they need to know? What did we learn so the next night
goes more smoothly? Those questions do not stop when commissioning ends. They follow the
observatory into full operations, and the platform has to follow as well.